\documentclass{ximera}

\input{../../../preamble.tex}


\definecolor{fvdnavy}{HTML}{071D43}
\definecolor{fvdcyan}{HTML}{20BFC6}
\definecolor{fvdcyanlight}{HTML}{EFFCFB}
\definecolor{fvdmagenta}{HTML}{F10D69}
\definecolor{fvdmagentalight}{HTML}{FFF1F7}
\definecolor{fvdtext}{HTML}{163044}





%=================================================
% Pedagogische omgevingen
% PDF: tcolorbox
% HTML: learning-box
%=================================================

\newenvironment{voorkennis}
{%
    \ifdefined\HCode
        \HCode{%
<section class="learning-box learning-box--prior">
<div class="learning-box__header">
<span class="learning-box__icon" aria-hidden="true"></span>
<span class="learning-box__title">Voorkennis</span>
</div>
<div class="learning-box__content">
        }%
        \begin{itemize}
    \else
       \begin{tcolorbox}[
    enhanced,
    breakable,
    colback=fvdcyanlight,
    colframe=fvdcyan,
    coltext=fvdtext,
    boxrule=0.5pt,
    leftrule=4pt,
    arc=4mm,
    outer arc=4mm,
    left=7mm,
    right=6mm,
    top=5mm,
    bottom=5mm,
    before skip=6mm,
    after skip=6mm,
    title=Voorkennis,
    coltitle=fvdcyan,
    fonttitle=\bfseries\Large,
    attach title to upper={\par\smallskip},
    boxed title style={
        empty,
        left=0mm,
        right=0mm,
        top=0mm,
        bottom=0mm
    },
    overlay={
        \draw[
            fvdcyan,
            line width=0.5pt
        ]
        ([xshift=42mm,yshift=-13mm]frame.north west)
        --
        ([xshift=-7mm,yshift=-13mm]frame.north east);
    }
]
        \begin{itemize}
    \fi
}
{%
    \end{itemize}
    \ifdefined\HCode
        \HCode{%
</div>
</section>
        }%
    \else
        \end{tcolorbox}
    \fi
}


\newenvironment{leerdoelen}
{%
    \ifdefined\HCode
        \HCode{%
<section class="learning-box learning-box--goals">
<div class="learning-box__header">
<span class="learning-box__icon" aria-hidden="true"></span>
<span class="learning-box__title">Leerdoelen</span>
</div>
<div class="learning-box__content">
        }%
        \begin{itemize}
    \else
\begin{tcolorbox}[
    enhanced,
    breakable,
    colback=fvdmagentalight,
    colframe=fvdmagenta,
    coltext=fvdtext,
    boxrule=0.5pt,
    leftrule=4pt,
    arc=4mm,
    outer arc=4mm,
    left=7mm,
    right=6mm,
    top=5mm,
    bottom=5mm,
    before skip=6mm,
    after skip=6mm,
    title=Leerdoelen,
    coltitle=fvdmagenta,
    fonttitle=\bfseries\Large,
    attach title to upper={\par\smallskip},
    boxed title style={
        empty,
        left=0mm,
        right=0mm,
        top=0mm,
        bottom=0mm
    },
    overlay={
        \draw[
            fvdmagenta,
            line width=0.5pt
        ]
        ([xshift=42mm,yshift=-13mm]frame.north west)
        --
        ([xshift=-7mm,yshift=-13mm]frame.north east);
    }
]
        \begin{itemize}
    \fi
}
{%
    \end{itemize}
    \ifdefined\HCode
        \HCode{%
</div>
</section>
        }%
    \else
        \end{tcolorbox}
    \fi
}


\addPrintStyle{../../..}

\title{Arbeid en energie}
\author{Ingmar Herreman}

\begin{abstract}
In de vorige hoofdstukken beschreven we bewegingen en onderzochten we hoe
krachten die bewegingen veranderen. In dit hoofdstuk bekijken we dezelfde
fysica vanuit een andere invalshoek: energie.

We starten met het begrip arbeid. Eerst herhalen we arbeid bij een constante
kracht en onderzoeken we welke component van een kracht werkelijk bijdraagt
aan een verplaatsing. Daarna lezen we arbeid rechtstreeks af uit een
kracht-plaatsgrafiek. Voor eenvoudige grafieken gebruiken we oppervlakten van
rechthoeken, driehoeken en trapezia. Wanneer de kracht continu verandert,
wordt diezelfde oppervlakte beschreven door een bepaalde integraal.

Vervolgens verbinden we arbeid met kinetische energie via het
arbeid-energietheorema. We onderzoeken conservatieve en niet-conservatieve
krachten en bouwen van daaruit de begrippen potentiële en mechanische energie
verder op. Zwaarte-energie nabij het aardoppervlak en elastische energie
worden opnieuw bekeken vanuit arbeid.

Ten slotte keren we terug naar gravitatie op astronomische schaal. Daar is
$mgh$ niet langer voldoende en gebruiken we de algemene gravitatie-energie.
Zo kunnen we satellieten nu niet alleen vanuit krachten, maar ook vanuit
energie begrijpen.

Doorheen het hoofdstuk staat één idee centraal:
\[
\boxed{\text{arbeid is een manier waarop energie wordt overgedragen of omgezet}.}
\]
Grafieken, oppervlakten en integralen vormen daarbij de brug tussen wiskunde
en fysica.
\end{abstract}

\begin{document}

% FVD_CHAPTER_DATA_START
% Automatisch gegenereerd uit index.tex.
% Niet handmatig aanpassen.
\fvdchapterdata{1}{Beweging en kracht — Van beweging beschrijven tot energie begrijpen}{8}
% FVD_CHAPTER_DATA_END



\maketitle

\begin{voorkennis}
  \item Je kan krachten tekenen, ontbinden en de resulterende kracht bepalen.
  \item Je kan het scalair product interpreteren via de hoek tussen vectoren.
  \item Je kent kinetische, gravitationele en elastische energie uit eerdere leerstof.
  \item Je kan eenvoudige energieomzettingen beschrijven.
  \item Je kan oppervlakten van rechthoeken, driehoeken en trapezia berekenen.
  \item Je kent de bepaalde integraal als georiënteerde oppervlakte.
  \item Je kan Newton II toepassen en snelheidsveranderingen analyseren.
\end{voorkennis}

\begin{leerdoelen}
  \item Je kan arbeid fysisch interpreteren als energieoverdracht of energieomzetting.
  \item Je kan arbeid van een constante kracht berekenen met
        $W=F\Delta r\cos\theta$.
  \item Je kan positieve, negatieve en nul-arbeid fysisch interpreteren.
  \item Je kan arbeid bepalen uit een $F_x(x)$-grafiek met de oppervlaktemethode.
  \item Je kan arbeid bij een niet-constante kracht bepalen met een bepaalde integraal.
  \item Je kan het arbeid-energietheorema gebruiken en interpreteren.
  \item Je kan conservatieve en niet-conservatieve krachten onderscheiden.
  \item Je kan potentiële energie verbinden met arbeid van een conservatieve kracht.
  \item Je kan energieomzettingen analyseren met behoud van mechanische energie
        wanneer het model dit toelaat.
  \item Je kan energiedissipatie analyseren wanneer niet-conservatieve krachten werken.
  \item Je kan arbeid- en energieproblemen systematisch mathematiseren,
        oplossen en fysisch interpreteren.
\end{leerdoelen}


% ============================================================
\FVDsection{Van kracht naar energie}
% ============================================================

In de mechanica kunnen we een beweging op verschillende manieren onderzoeken.

Met Newton II kijken we naar de krachten:

\[
\sum\vec F=m\vec a.
\]

Daarmee kunnen we bepalen hoe de snelheid op elk moment verandert.

Maar soms interesseert ons vooral het verschil tussen een begin- en
eindtoestand. Hoe snel rijdt een wagen na het remmen? Hoe hoog geraakt een
voorwerp? Hoeveel energie moet een motor leveren? Hoeveel energie gaat door
wrijving over in interne energie?

Dan is een energiebeschrijving vaak efficiënter.

\begin{center}
\[
\boxed{
\text{kracht over een verplaatsing}
\longrightarrow
\text{arbeid}
\longrightarrow
\text{verandering van energie}
}
\]
\end{center}

\begin{oefening}[Arbeid in het dagelijks leven en in fysica]{\oefB}

Beoordeel telkens of de genoemde kracht mechanische arbeid verricht op het
beschouwde voorwerp.

\begin{enumerate}
  \item Je duwt een slede vooruit en de slede verplaatst zich.
  \item Je duwt hard tegen een kast, maar de kast blijft staan.
  \item Je houdt een zware halter onbeweeglijk boven je hoofd.
  \item Een bloempot valt naar beneden.
  \item Je rekt een veer uit.
  \item De normaalkracht werkt op een kist die horizontaal over een vlak schuift.
\end{enumerate}

Motiveer telkens vanuit kracht én verplaatsing.

\end{oefening}


% ============================================================
\FVDsection{Arbeid van een constante kracht}
% ============================================================

Beschouw eerst een constante kracht die evenwijdig en gelijkgericht is met
de verplaatsing. Dan geldt:

\[
\boxed{W=F\Delta r.}
\]

De SI-eenheid van arbeid is de joule:

\[
\boxed{1\,\mathrm J=1\,\mathrm{N\,m}.}
\]

Arbeid is geen vector maar een scalaire grootheid.

\begin{opmerking}
Dat een kracht op een voorwerp werkt, betekent niet automatisch dat die kracht
arbeid verricht. Er moet een verplaatsing zijn én de kracht moet een component
hebben in de richting van die verplaatsing.
\end{opmerking}

\begin{oefening}[Constante kracht]{\oefB}

Een karretje wordt over $6{,}0\,\mathrm m$ voortgetrokken door een constante
horizontale kracht van $35\,\mathrm N$.

\begin{question}
Bereken de arbeid van de trekkracht.
\end{question}

\begin{question}
Hoe verandert de arbeid als dezelfde kracht over een dubbele afstand werkt?
\end{question}

\begin{question}
Welke energie-eenheid is equivalent met $\mathrm{N\,m}$?
\end{question}

\end{oefening}


% ============================================================
\FVDsection{Wanneer kracht en verplaatsing niet evenwijdig zijn}
% ============================================================

Een trekkracht staat vaak onder een hoek met de verplaatsing. Denk aan een
reiskoffer of slee die met een schuin touw wordt voortgetrokken.

Alleen de component van de kracht in de richting van de verplaatsing draagt
bij tot de arbeid.

Als $\theta$ de hoek is tussen $\vec F$ en $\Delta\vec r$, dan:

\[
F_\parallel=F\cos\theta.
\]

Daarom:

\[
W=F_\parallel\Delta r
\]

en dus:

\[
\boxed{
W=F\Delta r\cos\theta.
}
\]

Dit is precies het scalair product:

\[
\boxed{
W=\vec F\cdot\Delta\vec r.
}
\]

De formule bevat dus zowel de grootte van de kracht en verplaatsing als hun
onderlinge richting.


\FVDsubsection{Het teken van arbeid}

Uit de cosinus volgt onmiddellijk:

\[
0^\circ\leq\theta<90^\circ
\quad\Rightarrow\quad
W>0,
\]

\[
\theta=90^\circ
\quad\Rightarrow\quad
W=0,
\]

\[
90^\circ<\theta\leq180^\circ
\quad\Rightarrow\quad
W<0.
\]

Positieve arbeid betekent dat de betrokken kracht energie toevoegt aan de
bewegingsvorm die we onderzoeken. Negatieve arbeid betekent dat de kracht
daar energie aan onttrekt.

\begin{opmerking}
``Negatieve arbeid'' betekent niet dat er ``negatieve energie'' bestaat.
Het teken geeft de richting van de energieoverdracht aan binnen het gekozen
model.
\end{opmerking}


\begin{oefening}[Kracht onder een hoek]{\oefB\ESS}

Een persoon trekt een koffer met een constante kracht van $200\,\mathrm N$
over $100\,\mathrm m$. De trekkracht maakt een hoek van $30^\circ$ met de
horizontale verplaatsing.

\begin{question}
Bereken de component van de trekkracht in de bewegingsrichting.
\end{question}

\begin{question}
Bereken de arbeid van de trekkracht.
\end{question}

\begin{question}
Welke component van de trekkracht verricht in dit model geen arbeid?
Waarom?
\end{question}

\end{oefening}


\begin{oefening}[Positief, negatief of nul?]{\oefU\ESS}

Een kist schuift naar rechts over een horizontale vloer. Op de kist werken
zwaartekracht, normaalkracht, een trekkracht schuin naar rechtsboven en
wrijvingskracht.

\begin{question}
Maak een vrijlichaamsdiagram.
\end{question}

\begin{question}
Bepaal voor elke kracht of haar arbeid positief, negatief of nul is.
\end{question}

\begin{question}
Kan een kracht groot zijn en toch geen arbeid verrichten? Gebruik deze
situatie om je antwoord te motiveren.
\end{question}

\begin{question}
Waarom is het onvoldoende om alleen naar de richting van de kracht te kijken?
\end{question}

\end{oefening}


% ============================================================
\FVDsection{Arbeid uit een kracht-plaatsgrafiek}
% ============================================================

De formule

\[
W=F_x\Delta x
\]

heeft een belangrijke grafische betekenis.

Stel dat een constante krachtcomponent $F_x$ wordt uitgezet als functie van
de plaats $x$. Tussen $x_1$ en $x_2$ ontstaat onder de grafiek een rechthoek.

De oppervlakte is:

\[
F_x(x_2-x_1).
\]

Dat is precies de arbeid:

\[
\boxed{
W=\text{georiënteerde oppervlakte tussen de }F_x(x)\text{-grafiek en de }x\text{-as}.
}
\]

Deze interpretatie is veel algemener dan de formule voor een constante kracht.


\FVDsubsection{De oppervlaktemethode}

Voor eenvoudige, stukgewijs lineaire grafieken hoeven we niet onmiddellijk
een integraal te berekenen. We verdelen het gebied in gekende meetkundige
figuren.

Voor een rechthoek:

\[
W=bh.
\]

Voor een driehoek:

\[
W=\frac12bh.
\]

Voor een trapezium:

\[
W=\frac{h_1+h_2}{2}b.
\]

Hierbij stelt de horizontale afmeting een verplaatsing voor en de verticale
afmeting een krachtcomponent.

Controle van de eenheid:

\[
\mathrm{N\,m}=\mathrm J.
\]


\begin{oefening}[Oppervlaktemethode: rechthoek en driehoek]{\oefB\ESS}

De krachtcomponent $F_x$ op een voorwerp verloopt als volgt:

\begin{itemize}
  \item van $x=0$ tot $x=4\,\mathrm m$: $F_x=6\,\mathrm N$;
  \item van $x=4$ tot $x=8\,\mathrm m$: $F_x$ daalt lineair van
        $6\,\mathrm N$ tot $0\,\mathrm N$.
\end{itemize}

\begin{question}
Schets de $F_x(x)$-grafiek.
\end{question}

\begin{question}
Bereken de arbeid van $0$ tot $4\,\mathrm m$ als oppervlakte van een rechthoek.
\end{question}

\begin{question}
Bereken de arbeid van $4$ tot $8\,\mathrm m$ als oppervlakte van een driehoek.
\end{question}

\begin{question}
Bepaal de totale arbeid.
\end{question}

\end{oefening}


\FVDsubsection{Oppervlakte onder de as}

Wanneer $F_x<0$ terwijl de positieve $x$-richting als bewegingsrichting is
gekozen, ligt de grafiek onder de $x$-as. Die oppervlakte telt negatief mee.

We spreken daarom nauwkeuriger van een \textbf{georiënteerde oppervlakte}.

\begin{oefening}[Georiënteerde oppervlakte]{\oefU\ESS}

Een krachtcomponent is:

\[
F_x=
\begin{cases}
4\,\mathrm N, & 0\leq x<3\,\mathrm m,\\
-2\,\mathrm N, & 3\leq x\leq7\,\mathrm m.
\end{cases}
\]

\begin{question}
Schets de grafiek.
\end{question}

\begin{question}
Bereken de positieve oppervlakte.
\end{question}

\begin{question}
Bereken de negatieve bijdrage.
\end{question}

\begin{question}
Bepaal de netto arbeid.
\end{question}

\begin{question}
Leg uit waarom ``totale oppervlakte'' zonder teken hier tot een fout antwoord
kan leiden.
\end{question}

\end{oefening}


% ============================================================
\FVDsection{Niet-constante kracht: van oppervlakte naar integraal}
% ============================================================

In veel realistische situaties verandert een kracht voortdurend met de plaats.
Dan bestaat de oppervlakte onder de $F_x(x)$-grafiek niet uit één eenvoudige
meetkundige figuur.

De bepaalde integraal veralgemeent precies de oppervlaktemethode:

\[
\boxed{
W_{x_1\rightarrow x_2}
=
\int_{x_1}^{x_2}F_x(x)\,dx.
}
\]

De integraal is hier geen los wiskundig trucje. Ze drukt hetzelfde fysische
idee uit als de oppervlakteberekeningen uit de vorige paragraaf.

\begin{center}
\[
\boxed{
\text{oppervlaktemethode}
\longrightarrow
\text{bepaalde integraal}
}
\]
\end{center}

Als $F_x$ constant is:

\[
W=\int_{x_1}^{x_2}F_x\,dx
=F_x(x_2-x_1),
\]

zodat de bekende formule terug verschijnt.


\begin{oefening}[Een lineair veranderende kracht]{\oefU\ESS}

Een krachtcomponent wordt beschreven door:

\[
F_x(x)=3x+2
\]

met $F_x$ in newton en $x$ in meter. Het voorwerp beweegt van $x=0$ tot
$x=4\,\mathrm m$.

\begin{question}
Schets de grafiek.
\end{question}

\begin{question}
Bereken de arbeid eerst met de oppervlaktemethode.
\end{question}

\begin{question}
Bereken dezelfde arbeid met:
\[
W=\int_0^4(3x+2)\,dx.
\]
\end{question}

\begin{question}
Vergelijk beide resultaten en verklaar waarom ze gelijk moeten zijn.
\end{question}

\end{oefening}


\begin{oefening}[Een niet-lineaire kracht]{\oefU\ESS}

Voor een krachtcomponent geldt:

\[
F_x(x)=2x^2+1.
\]

Een voorwerp beweegt van $x=1\,\mathrm m$ naar $x=3\,\mathrm m$.

\begin{question}
Waarom is een eenvoudige driehoek- of trapeziumformule hier niet exact?
\end{question}

\begin{question}
Stel de bepaalde integraal voor de arbeid op.
\end{question}

\begin{question}
Bereken de arbeid.
\end{question}

\begin{question}
Controleer de eenheid van je resultaat.
\end{question}

\end{oefening}


% ============================================================
\FVDsection{De veerkracht: een belangrijke variabele kracht}
% ============================================================

Voor een ideale veer geldt de wet van Hooke:

\[
\boxed{
F_{\mathrm{veer}}=-kx.
}
\]

Hierbij is $x$ de uitwijking ten opzichte van de evenwichts- of ongerekte
toestand. De constante $k$ is de veerconstante.

Het minteken toont dat de veerkracht een terugroepkracht is: ze is tegengesteld
gericht aan de uitwijking.

Let goed op de betekenis van $x$. Als een veer een rustlengte $L_0$ heeft en
een lengte $L$, dan is:

\[
x=L-L_0.
\]

We integreren dus over de \emph{uitwijking}, niet automatisch over de absolute
lengte van de veer.


\FVDsubsection{Arbeid van de veerkracht}

De arbeid van de veerkracht bij een verandering van $x_1$ naar $x_2$ is:

\[
W_{\mathrm{veer}}
=
\int_{x_1}^{x_2}-kx\,dx.
\]

Dus:

\[
\boxed{
W_{\mathrm{veer}}
=
-\frac12k(x_2^2-x_1^2).
}
\]

Bij het uitrekken van $0$ tot $x>0$:

\[
W_{\mathrm{veer}}=-\frac12kx^2.
\]

De externe kracht die de veer langzaam uitrekt, verricht in het ideale
quasistatische geval positieve arbeid:

\[
W_{\mathrm{ext}}=+\frac12kx^2.
\]


\begin{oefening}[Veer: grafiek, oppervlakte en integraal]{\oefU\ESS}

Een ideale veer heeft:

\[
k=240\,\mathrm{\frac Nm}.
\]

Ze wordt vanuit haar ongerekte toestand uitgerekt tot:

\[
x=0{,}050\,\mathrm m.
\]

\begin{question}
Schets $F_{\mathrm{veer}}(x)$.
\end{question}

\begin{question}
Bepaal de arbeid van de veerkracht met de oppervlaktemethode.
\end{question}

\begin{question}
Controleer het resultaat met een bepaalde integraal.
\end{question}

\begin{question}
Bepaal de arbeid van een externe kracht wanneer de veer zeer langzaam wordt
uitgerekt.
\end{question}

\begin{question}
Waarom hebben beide arbeidsgrootheden een tegengesteld teken?
\end{question}

\end{oefening}


% ============================================================
\FVDsection{Kinetische energie}
% ============================================================

Een bewegend voorwerp bezit kinetische energie:

\[
\boxed{
E_k=\frac12mv^2.
}
\]

De kinetische energie hangt af van de massa en van het kwadraat van de
snelheid.

Daarom:

\[
v\rightarrow2v
\quad\Rightarrow\quad
E_k\rightarrow4E_k.
\]

En:

\[
m\rightarrow2m
\quad\Rightarrow\quad
E_k\rightarrow2E_k.
\]

Omdat $v^2\geq0$ is:

\[
E_k\geq0.
\]


\begin{oefening}[Snelheid telt kwadratisch]{\oefB}

Een wagen rijdt eerst met snelheid $v$ en daarna met snelheid $2v$. De massa
blijft constant.

\begin{question}
Met welke factor verandert de kinetische energie?
\end{question}

\begin{question}
Waarom is de uitspraak ``dubbele snelheid betekent dubbele kinetische
energie'' fout?
\end{question}

\begin{question}
Wat betekent dit kwalitatief voor de energie die bij een noodstop moet
worden omgezet?
\end{question}

\end{oefening}


% ============================================================
\FVDsection{Het arbeid-energietheorema}
% ============================================================

De arbeid van alle krachten samen is gelijk aan de verandering van de
kinetische energie:

\[
\boxed{
W_{\mathrm{res}}
=
\Delta E_k.
}
\]

Dus:

\[
\boxed{
W_{\mathrm{res}}
=
\frac12mv_e^2-\frac12mv_b^2.
}
\]

Dit is het \textbf{arbeid-energietheorema}.

Het woord ``res'' is belangrijk. We mogen ofwel de arbeid van de resulterende
kracht bepalen, ofwel de arbeid van alle afzonderlijke krachten optellen:

\[
W_{\mathrm{res}}
=
\sum_i W_i.
\]

Als:

\[
W_{\mathrm{res}}>0,
\]

dan neemt de kinetische energie toe.

Als:

\[
W_{\mathrm{res}}<0,
\]

dan neemt de kinetische energie af.

Als:

\[
W_{\mathrm{res}}=0,
\]

dan blijft de kinetische energie constant.


\FVDsubsection{Waarom werkt het?}

In één dimensie:

\[
F_{\mathrm{res}}=ma.
\]

Met:

\[
a=\frac{dv}{dt}
\qquad\text{en}\qquad
v=\frac{dx}{dt}
\]

kunnen we schrijven:

\[
a
=
\frac{dv}{dt}
=
\frac{dv}{dx}\frac{dx}{dt}
=
v\frac{dv}{dx}.
\]

Daarom:

\[
F_{\mathrm{res}}
=
mv\frac{dv}{dx}.
\]

Vermenigvuldigen met $dx$:

\[
F_{\mathrm{res}}\,dx=m v\,dv.
\]

Integreren tussen begin- en eindtoestand:

\[
\int_{x_b}^{x_e}F_{\mathrm{res}}\,dx
=
\int_{v_b}^{v_e}mv\,dv.
\]

Dus:

\[
W_{\mathrm{res}}
=
\left[\frac12mv^2\right]_{v_b}^{v_e}
\]

en:

\[
\boxed{
W_{\mathrm{res}}=\Delta E_k.
}
\]

Deze afleiding toont een diepe verbinding tussen Newtons bewegingswet en de
energiebeschrijving.


\begin{oefening}[Versnellen via arbeid]{\oefB\ESS}

Een auto met massa $1000\,\mathrm{kg}$ rijdt aanvankelijk
$3{,}0\,\mathrm{m/s}$. Een constante resulterende kracht van
$4000\,\mathrm N$ werkt over een afstand van $5{,}0\,\mathrm m$ in de
bewegingsrichting.

\begin{question}
Bereken de arbeid van de resulterende kracht.
\end{question}

\begin{question}
Gebruik het arbeid-energietheorema om de eindsnelheid te bepalen.
\end{question}

\begin{question}
Welke alternatieve methode met Newton II en kinematica zou eveneens mogelijk
zijn?
\end{question}

\begin{question}
Welke methode vind je hier het kortst? Motiveer.
\end{question}

\end{oefening}


\begin{oefening}[Remweg]{\oefU\ESS}

Een voertuig wordt geremd door een ongeveer constante resulterende remkracht.

\begin{question}
Gebruik het arbeid-energietheorema om aan te tonen dat, bij dezelfde massa en
dezelfde remkracht,
\[
d_{\mathrm{rem}}\sim v_0^2.
\]
\end{question}

\begin{question}
Wat gebeurt er met de remweg als de beginsnelheid verdubbelt?
\end{question}

\begin{question}
Waarom is dit resultaat relevant voor verkeersveiligheid?
\end{question}

\begin{question}
Welke vereenvoudigingen bevat het model?
\end{question}

\end{oefening}


% ============================================================
\FVDsection{Conservatieve krachten}
% ============================================================

Bij sommige krachten hangt de verrichte arbeid tussen twee punten alleen af
van begin- en eindpunt, niet van het gevolgde pad.

Zo'n kracht noemen we \textbf{conservatief}.

Voor een conservatieve kracht kunnen we een potentiële energie definiëren:

\[
\boxed{
W_{\mathrm{cons}}=-\Delta E_p.
}
\]

Dus:

\[
W_{\mathrm{cons}}
=
E_{p,b}-E_{p,e}.
\]

Voorbeelden die we in deze cursus gebruiken zijn:

\begin{itemize}
  \item gravitatiekracht;
  \item zwaartekracht in de nabij-aardebenadering;
  \item ideale veerkracht.
\end{itemize}


\FVDsubsection{Niet-conservatieve krachten}

Bij niet-conservatieve krachten hangt de arbeid in het algemeen wel af van het
gevolgde pad of treedt overdracht naar andere energievormen op.

Wrijving is het belangrijkste mechanische voorbeeld in deze cursus.

Mechanische energie hoeft dan niet behouden te blijven.

Dat betekent niet dat energie ``verdwijnt''. Ze wordt bijvoorbeeld omgezet in
interne energie van voorwerp en omgeving.


\begin{oefening}[Conservatief of niet?]{\oefU\ESS}

Beschouw gravitatiekracht, ideale veerkracht en kinetische wrijvingskracht.

\begin{question}
Welke krachten modelleren we als conservatief?
\end{question}

\begin{question}
Voor welke krachten kunnen we in deze context een potentiële energie
definiëren?
\end{question}

\begin{question}
Waarom leidt wrijving vaak tot een afname van de mechanische energie?
\end{question}

\begin{question}
Is de totale energie daarom niet behouden? Leg zorgvuldig uit.
\end{question}

\end{oefening}


% ============================================================
\FVDsection{Zwaarte-energie nabij het aardoppervlak}
% ============================================================

Nabij het aardoppervlak behandelen we het gravitatieveld vaak als homogeen:

\[
g\approx\text{constant}.
\]

Kies de positieve $y$-richting omhoog. De zwaartekracht is dan:

\[
F_{g,y}=-mg.
\]

Voor een verplaatsing van $y_1$ naar $y_2$:

\[
W_g
=
\int_{y_1}^{y_2}-mg\,dy
=
-mg(y_2-y_1).
\]

Omdat:

\[
W_g=-\Delta E_{p,g},
\]

volgt:

\[
\Delta E_{p,g}=mg(y_2-y_1).
\]

Als we een referentieniveau $h=0$ kiezen:

\[
\boxed{
E_{p,g}=mgh.
}
\]

De absolute waarde hangt af van de gekozen nulhoogte; energieverschillen zijn
fysisch relevant.


\begin{oefening}[De nulhoogte kiezen]{\oefB}

Een boek met massa $1{,}5\,\mathrm{kg}$ ligt op een tafel van
$0{,}80\,\mathrm m$ hoog.

\begin{question}
Bereken de potentiële zwaarte-energie als de vloer het nulniveau is.
\end{question}

\begin{question}
Welke waarde krijg je als het tafelblad als nulniveau wordt gekozen?
\end{question}

\begin{question}
Verandert de arbeid van de zwaartekracht wanneer het boek van de tafel op de
vloer valt door die andere keuze? Verklaar.
\end{question}

\end{oefening}


% ============================================================
\FVDsection{Elastische potentiële energie}
% ============================================================

We vonden voor de arbeid van een ideale veerkracht:

\[
W_{\mathrm{veer}}
=
-\frac12k(x_2^2-x_1^2).
\]

Voor een conservatieve kracht:

\[
W_{\mathrm{veer}}=-\Delta E_{p,\mathrm{el}}.
\]

Dus:

\[
\Delta E_{p,\mathrm{el}}
=
\frac12k(x_2^2-x_1^2).
\]

Met de ongerekte toestand als referentie:

\[
\boxed{
E_{p,\mathrm{el}}=\frac12kx^2.
}
\]

Dit resultaat is dus niet zomaar een formule om te onthouden: het volgt uit
de oppervlakte onder de kracht-plaatsgrafiek.


\begin{oefening}[Van $F(x)$ naar energie]{\oefU\ESS}

Een ideale veer heeft veerconstante $k$.

\begin{question}
Schets de grootte van de externe kracht die nodig is om de veer zeer langzaam
van $x=0$ tot $x=X$ uit te rekken.
\end{question}

\begin{question}
Gebruik de oppervlakte van de driehoek om de externe arbeid te bepalen.
\end{question}

\begin{question}
Toon zo aan dat:
\[
E_{p,\mathrm{el}}=\frac12kX^2.
\]
\end{question}

\begin{question}
Welke wiskundige methode geeft hetzelfde resultaat voor een willekeurige
krachtfunctie?
\end{question}

\end{oefening}


% ============================================================
\FVDsection{Mechanische energie}
% ============================================================

De mechanische energie van een systeem is de som van kinetische en relevante
potentiële energieën:

\[
\boxed{
E_m=E_k+E_p.
}
\]

Wanneer alleen conservatieve krachten arbeid verrichten:

\[
\boxed{
E_{m,b}=E_{m,e}.
}
\]

Of:

\[
\boxed{
E_k+E_p=\text{constant}.
}
\]

Energie kan voortdurend van vorm veranderen terwijl de totale mechanische
energie constant blijft.


\begin{oefening}[Vallen met energie]{\oefB\ESS}

Een voorwerp wordt vanuit rust losgelaten op hoogte $h$. Verwaarloos
luchtweerstand.

\begin{question}
Kies het grondniveau als nulniveau en schrijf de mechanische energie bij het
loslaten.
\end{question}

\begin{question}
Schrijf de mechanische energie vlak voor het bereiken van de grond.
\end{question}

\begin{question}
Leid daaruit af:
\[
v=\sqrt{2gh}.
\]
\end{question}

\begin{question}
Vergelijk deze methode met de kinematische aanpak uit het hoofdstuk vrije val.
\end{question}

\end{oefening}


\begin{oefening}[Veer lanceert een karretje]{\oefU\ESS}

Een karretje met massa $m$ staat op een horizontaal wrijvingsloos spoor tegen
een ingedrukte veer met veerconstante $k$. De indrukking is $x$.

\begin{question}
Welke energievorm is aanvankelijk aanwezig?
\end{question}

\begin{question}
Welke energievorm bezit het karretje wanneer de veer haar evenwichtsstand
bereikt?
\end{question}

\begin{question}
Leid een formule af voor de snelheid:
\[
v=x\sqrt{\frac{k}{m}}.
\]
\end{question}

\begin{question}
Wat gebeurt er met $v$ als de indrukking verdubbelt?
\end{question}

\end{oefening}


% ============================================================
\FVDsection{Niet-conservatieve arbeid en energiedissipatie}
% ============================================================

Als niet-conservatieve krachten arbeid verrichten, is de mechanische energie
in het algemeen niet constant.

We kunnen schrijven:

\[
\boxed{
W_{\mathrm{nc}}=\Delta E_m
}
\]

wanneer $W_{\mathrm{nc}}$ de arbeid van de niet-conservatieve krachten op het
gekozen mechanische systeem voorstelt.

Bij kinetische wrijving is die arbeid vaak negatief:

\[
W_w=-F_w d
\]

voor een constante wrijvingskracht tegengesteld aan de beweging.

De ``verdwenen'' mechanische energie verschijnt als andere energievormen,
bijvoorbeeld interne energie.


\begin{oefening}[Glijden met wrijving]{\oefU\ESS}

Een blok met massa $m$ vertrekt vanuit rust van hoogte $h$ en glijdt naar
beneden. Door wrijving wordt een hoeveelheid mechanische energie
$E_{\mathrm{diss}}$ gedissipeerd.

\begin{question}
Stel een energiebalans op.
\end{question}

\begin{question}
Leid een uitdrukking af voor de snelheid onderaan.
\end{question}

\begin{question}
Vergelijk met het wrijvingsloze geval.
\end{question}

\begin{question}
Waarom is het fysisch misleidend om te zeggen dat de energie door wrijving
``verloren'' is?
\end{question}

\end{oefening}


% ============================================================
\FVDsection{Wanneer gebruik je Newton en wanneer energie?}
% ============================================================

Beide methoden beschrijven dezelfde fysica, maar beantwoorden vaak andere
vragen efficiënt.

Gebruik een krachtenaanpak wanneer je bijvoorbeeld de versnelling, richting
van de versnelling of tijdsontwikkeling nodig hebt.

Gebruik een energieaanpak wanneer vooral begin- en eindtoestand belangrijk
zijn en het gevolgde tijdsverloop niet nodig is.

Soms moet je beide combineren.

\begin{oefening}[Kies een methode]{\oefU\ESS}

Geef voor elke vraag aan of een aanpak via Newton/kinematica, via energie, of
een combinatie het meest efficiënt lijkt. Motiveer.

\begin{enumerate}
  \item Bepaal de snelheid van een vallend voorwerp na een hoogteverschil $h$.
  \item Bepaal de normaalkracht op een auto in het laagste punt van een
        verticale cirkelbaan.
  \item Bepaal hoe ver een blok met beginsnelheid $v_0$ over een ruw vlak
        schuift voor het stopt.
  \item Bepaal de versnelling van een blok op een hellend vlak.
  \item Bepaal de maximale hoogte van een karretje op een wrijvingsloze baan.
\end{enumerate}

\end{oefening}


% ============================================================
\FVDsection{Algemene gravitatie-energie}
% ============================================================

In H7 gebruikten we:

\[
F_g=G\frac{Mm}{r^2}.
\]

Voor satellieten verandert $r$ aanzienlijk. De nabij-aardeformule

\[
E_{p,g}=mgh
\]

is dan geen geschikte algemene beschrijving.

We kiezen de gravitatiepotentiële energie gelijk aan nul op oneindig grote
afstand:

\[
E_{p,g}(\infty)=0.
\]

Voor de gravitatiekracht volgt dan:

\[
\boxed{
E_{p,g}(r)=-\frac{GMm}{r}.
}
\]

De negatieve waarde heeft een duidelijke betekenis: een massa op eindige
afstand bevindt zich in een gebonden toestand ten opzichte van het gekozen
nulniveau op oneindig.


\FVDsubsection{Waarom een minteken?}

Om een massa vanuit een eindige afstand $r$ tot oneindig te brengen zonder
eindsnelheid, moet van buitenaf positieve arbeid worden geleverd.

De potentiële energie stijgt daarbij van:

\[
-\frac{GMm}{r}
\]

naar:

\[
0.
\]

Daarom is de gravitatiepotentiële energie op eindige afstand negatief bij
deze conventie.


\begin{oefening}[Van dichtbij naar ver weg]{\oefU}

Een satelliet wordt van een cirkelbaan met straal $r_1$ naar een grotere
afstand $r_2>r_1$ gebracht.

\begin{question}
Vergelijk $E_{p,g}(r_1)$ en $E_{p,g}(r_2)$.
\end{question}

\begin{question}
Welke waarde ligt dichter bij nul?
\end{question}

\begin{question}
Waarom betekent ``grotere potentiële gravitatie-energie'' hier niet
noodzakelijk een positief getal?
\end{question}

\end{oefening}


% ============================================================
\FVDsection{Satellieten opnieuw bekeken: nu met energie}
% ============================================================

Voor een satelliet in een ideale cirkelbaan vonden we in H7:

\[
v^2=\frac{GM}{r}.
\]

De kinetische energie is dus:

\[
E_k
=
\frac12mv^2
=
\frac12m\frac{GM}{r}.
\]

Daarom:

\[
\boxed{
E_k=\frac{GMm}{2r}.
}
\]

De gravitatiepotentiële energie is:

\[
\boxed{
E_{p,g}=-\frac{GMm}{r}.
}
\]

De totale mechanische energie wordt:

\[
E_m=E_k+E_{p,g}.
\]

Dus:

\[
\boxed{
E_m=-\frac{GMm}{2r}.
}
\]

Voor een ideale cirkelbaan geldt bijgevolg:

\[
\boxed{
E_{p,g}=-2E_k
}
\]

en:

\[
\boxed{
E_m=-E_k.
}
\]


\begin{oefening}[Een hogere cirkelbaan]{\oefG}

Een satelliet wordt van een cirkelbaan met straal $r$ naar een cirkelbaan met
straal $4r$ gebracht.

\begin{question}
Met welke factor verandert de baansnelheid?
\end{question}

\begin{question}
Met welke factor verandert de kinetische energie?
\end{question}

\begin{question}
Vergelijk de gravitatiepotentiële energie van beide banen.
\end{question}

\begin{question}
Vergelijk de totale mechanische energie.
\end{question}

\begin{question}
Moet netto energie worden toegevoerd of afgevoerd om de satelliet naar de
hogere cirkelbaan te brengen? Motiveer vanuit $E_m$.
\end{question}

\end{oefening}


% ============================================================
\FVDsection{Gevorderd: ontsnappen aan een gravitatieveld}
% ============================================================

Hoe groot moet de beginsnelheid aan het oppervlak van een planeet minstens
zijn zodat een voorwerp in het ideale model oneindig ver kan raken zonder nog
kinetische energie over te houden?

We verwaarlozen atmosfeer, rotatie en andere hemellichamen.

Aan het oppervlak:

\[
E_{m,b}
=
\frac12mv_{\mathrm{esc}}^2-\frac{GMm}{R}.
\]

In de grenssituatie op oneindig:

\[
E_{m,e}=0.
\]

Behoud van mechanische energie geeft:

\[
\frac12mv_{\mathrm{esc}}^2-\frac{GMm}{R}=0.
\]

Dus:

\[
\boxed{
v_{\mathrm{esc}}
=
\sqrt{\frac{2GM}{R}}.
}
\]

De massa van het ontsnappende voorwerp valt opnieuw weg.

\begin{oefening}[Ontsnappingssnelheid]{\oefG}

\begin{question}
Leid de formule voor $v_{\mathrm{esc}}$ zelfstandig af vanuit energiebehoud.
\end{question}

\begin{question}
Vergelijk met de cirkelbaansnelheid vlak boven het oppervlak:
\[
v_{\mathrm{baan}}=\sqrt{\frac{GM}{R}}.
\]
\end{question}

\begin{question}
Toon aan:
\[
v_{\mathrm{esc}}=\sqrt2\,v_{\mathrm{baan}}.
\]
\end{question}

\begin{question}
Waarom is ontsnappingssnelheid in dit ideale model onafhankelijk van de massa
van het gelanceerde voorwerp?
\end{question}

\begin{question}
Welke belangrijke realistische effecten zijn niet in dit model opgenomen?
\end{question}

\end{oefening}


% ============================================================
\FVDsection{Een krachtgrafiek als volledig energieprobleem}
% ============================================================

Een $F_x(x)$-grafiek kan rechtstreeks vertellen hoe de kinetische energie
verandert:

\[
\Delta E_k
=
W_{\mathrm{res}}
=
\int_{x_1}^{x_2}F_{\mathrm{res},x}(x)\,dx.
\]

Dit maakt een grafiek tot veel meer dan een afbeelding van de kracht. De
georiënteerde oppervlakte bevat informatie over de verandering van de
beweging.


\begin{oefening}[Van grafiek naar snelheid]{\oefU\ESS}

Een voorwerp met massa $2{,}0\,\mathrm{kg}$ beweegt bij $x=0$ met
$3{,}0\,\mathrm{m/s}$. De resulterende kracht verloopt als volgt:

\begin{itemize}
  \item van $0$ tot $2\,\mathrm m$ stijgt $F_x$ lineair van $0$ tot
        $6\,\mathrm N$;
  \item van $2$ tot $5\,\mathrm m$ blijft $F_x=6\,\mathrm N$;
  \item van $5$ tot $7\,\mathrm m$ daalt $F_x$ lineair tot $0$.
\end{itemize}

\begin{question}
Schets de $F_x(x)$-grafiek.
\end{question}

\begin{question}
Bepaal met de oppervlaktemethode de totale arbeid.
\end{question}

\begin{question}
Bereken de kinetische energie bij $x=7\,\mathrm m$.
\end{question}

\begin{question}
Bereken de snelheid daar.
\end{question}

\begin{question}
Waarom hoefde je de tijdsduur van de beweging nergens te kennen?
\end{question}

\end{oefening}


\begin{oefening}[Een kracht die van teken verandert]{\oefG}

De resulterende kracht op een deeltje is:

\[
F_x(x)=6-2x
\]

voor $0\leq x\leq5\,\mathrm m$.

\begin{question}
Op welke plaats verandert de kracht van teken?
\end{question}

\begin{question}
Bepaal de arbeid van $x=0$ tot dat punt met een oppervlakte.
\end{question}

\begin{question}
Bepaal de arbeid over het volledige interval met een bepaalde integraal.
\end{question}

\begin{question}
Kan de kinetische energie op een deel van het traject toenemen en later
afnemen? Leg uit vanuit het teken van $F_x$.
\end{question}

\begin{question}
Waarom is het teken van de lokale kracht niet hetzelfde als het teken van de
totale arbeid over een volledig traject?
\end{question}

\end{oefening}


% ============================================================
\FVDsection{Problemen oplossen met energie}
% ============================================================

Gebruik bij een energieprobleem een vaste strategie.

\begin{enumerate}
  \item \textbf{Kies het systeem.} Welke objecten behoren tot het systeem?
  \item \textbf{Kies begin- en eindtoestand.}
  \item \textbf{Maak een schets.}
  \item \textbf{Inventariseer relevante energievormen.}
  \item \textbf{Bepaal of niet-conservatieve arbeid relevant is.}
  \item \textbf{Schrijf eerst symbolisch de energiebalans.}
  \item \textbf{Vul pas daarna waarden in.}
  \item \textbf{Controleer eenheden en grootteorde.}
  \item \textbf{Interpreteer het resultaat.}
\end{enumerate}

Een energietabel kan helpen:

\[
\begin{array}{c|c|c}
&\text{begin}&\text{einde}\\ \hline
E_k&\cdots&\cdots\\
E_{p,g}&\cdots&\cdots\\
E_{p,\mathrm{el}}&\cdots&\cdots
\end{array}
\]


\begin{oefening}[Achtbaan zonder motor]{\oefU\ESS}

Een karretje vertrekt vanuit rust op hoogte $h_1$ en bereikt later hoogte
$h_2<h_1$. Verwaarloos eerst wrijving.

\begin{question}
Maak een energietabel.
\end{question}

\begin{question}
Leid een formule af voor de snelheid op hoogte $h_2$.
\end{question}

\begin{question}
Welke massa-afhankelijkheid blijft in het eindresultaat over?
\end{question}

\begin{question}
Herhaal de energiebalans wanneer onderweg een hoeveelheid
$E_{\mathrm{diss}}$ wordt gedissipeerd.
\end{question}

\end{oefening}


% ============================================================
\FVDsection{STEM-context: regeneratief remmen}
% ============================================================

Bij klassieke wrijvingsremmen wordt een groot deel van de kinetische energie
omgezet in interne energie van remmen en omgeving.

Bij regeneratief remmen kan een elektromotor als generator werken en een deel
van de kinetische energie omzetten in elektrische energie die opnieuw wordt
opgeslagen.

Een eenvoudig energiemodel is:

\[
E_{\mathrm{terug}}=\eta\,\Delta E_k,
\]

waar $0\leq\eta\leq1$ het rendement van de terugwinning voorstelt.


\begin{oefening}[Energie terugwinnen]{\oefU}

Een elektrisch voertuig met massa $1600\,\mathrm{kg}$ vertraagt van
$20\,\mathrm{m/s}$ tot $5{,}0\,\mathrm{m/s}$. Van de afname van de kinetische
energie kan $65\%$ elektrisch worden teruggewonnen.

\begin{question}
Bereken de afname van de kinetische energie.
\end{question}

\begin{question}
Bereken hoeveel energie wordt teruggewonnen.
\end{question}

\begin{question}
Waar blijft de rest van de energie?
\end{question}

\begin{question}
Waarom kan een echt regeneratief remsysteem nooit alle kinetische energie
opnieuw als bruikbare elektrische energie opslaan?
\end{question}

\end{oefening}


% ============================================================
\FVDsection{STEM-context: botsing met een meteoor}
% ============================================================

Bij zeer hoge snelheden kan de kinetische energie enorm zijn door het
kwadratische snelheidsverband.

\begin{oefening}[Barringerkrater]{\oefB}

Een meteoor heeft vóór de inslag naar schatting een massa van
$65\,\mathrm{ton}$ en een snelheid van $15\,\mathrm{km/s}$.

\begin{question}
Zet beide grootheden om naar SI-eenheden.
\end{question}

\begin{question}
Bereken de kinetische energie vlak vóór de inslag.
\end{question}

\begin{question}
Waarom kan die energie na de inslag niet eenvoudig volledig als kinetische
energie van één object worden teruggevonden?
\end{question}

\begin{question}
Noem minstens drie mogelijke energievormen of processen waarin de
oorspronkelijke kinetische energie terechtkomt.
\end{question}

\end{oefening}


% ============================================================
\FVDsection{ICT: arbeid numeriek benaderen}
% ============================================================

Niet elke krachtfunctie heeft een eenvoudige primitieve die leerlingen meteen
kunnen gebruiken. Experimentele gegevens bestaan bovendien vaak uit een tabel
van gemeten waarden.

Dan kunnen we de arbeid numeriek benaderen.

Bij meetpunten:

\[
(x_0,F_0),(x_1,F_1),\ldots,(x_n,F_n)
\]

kan een trapeziumbenadering worden gebruikt:

\[
\boxed{
W\approx
\sum_{i=0}^{n-1}
\frac{F_i+F_{i+1}}{2}(x_{i+1}-x_i).
}
\]

Dit is de discrete versie van hetzelfde oppervlakte-idee.


\begin{oefening}[Arbeid uit meetgegevens]{\oefU}

Bij een trekproef worden kracht en verplaatsing gemeten:

\[
\begin{array}{c|ccccc}
x\,(\mathrm m)&0&0{,}10&0{,}20&0{,}30&0{,}40\\ \hline
F\,(\mathrm N)&0&18&35&49&60
\end{array}
\]

\begin{question}
Maak met ICT een spreidingsdiagram van $F$ tegen $x$.
\end{question}

\begin{question}
Benader de arbeid met trapezia.
\end{question}

\begin{question}
Welke fysische betekenis heeft de oppervlakte onder de gemeten grafiek?
\end{question}

\begin{question}
Waarom is een numerieke methode hier natuurlijker dan zomaar een
functievoorschrift te verzinnen?
\end{question}

\end{oefening}


% ============================================================
\FVDsection{Geïntegreerde analyse}
% ============================================================

\begin{oefening}[Kist op een helling]{\oefU\ESS}

Een kist glijdt vanuit rust een helling af over afstand $s$. De hellingshoek
is $\alpha$ en de kinetische wrijvingscoëfficiënt is $\mu_k$.

\begin{question}
Maak een vrijlichaamsdiagram.
\end{question}

\begin{question}
Bepaal de arbeid van de zwaartekracht.
\end{question}

\begin{question}
Bepaal de arbeid van de normaalkracht.
\end{question}

\begin{question}
Bepaal de arbeid van de wrijvingskracht.
\end{question}

\begin{question}
Gebruik het arbeid-energietheorema om een formule voor de eindsnelheid af te
leiden.
\end{question}

\begin{question}
Controleer of je resultaat overeenkomt met wat je via Newton II en kinematica
zou vinden.
\end{question}

\end{oefening}


\begin{oefening}[Welke informatie ontbreekt?]{\oefG}

Een grafiek toont de resulterende kracht $F_x(x)$ op een voorwerp. De
oppervlakte tussen $x_1$ en $x_2$ bedraagt $40\,\mathrm J$.

Een leerling besluit:

\begin{quote}
``De snelheid bij $x_2$ is dus $40\,\mathrm{m/s}$.''
\end{quote}

\begin{question}
Welke fysische grootheid stelt $40\,\mathrm J$ werkelijk voor?
\end{question}

\begin{question}
Welke bijkomende gegevens zijn minstens nodig om de eindsnelheid te bepalen?
\end{question}

\begin{question}
Stel een algemene formule voor $v_2$ op in functie van $W$, $m$ en $v_1$.
\end{question}

\begin{question}
Onder welke voorwaarde bestaat binnen dit klassieke model een reële
eindsnelheid?
\end{question}

\end{oefening}


% ============================================================
\FVDsection{Veelvoorkomende denkfouten}
% ============================================================

\begin{itemize}
  \item \textbf{``Er werkt een kracht, dus er is arbeid.''}\\
        Niet noodzakelijk. Zonder verplaatsing of bij een kracht loodrecht op
        de verplaatsing is de arbeid nul.

  \item \textbf{``Arbeid is altijd positief.''}\\
        Nee. Arbeid is een scalaire grootheid met teken.

  \item \textbf{``De oppervlakte onder een $F(x)$-grafiek is altijd positief.''}\\
        Nee. Gebieden onder de $x$-as leveren negatieve arbeid.

  \item \textbf{``$mgh$ is altijd de gravitatiepotentiële energie.''}\\
        Nee. $mgh$ is een lokale benadering voor een ongeveer homogeen
        gravitatieveld.

  \item \textbf{``Wrijving vernietigt energie.''}\\
        Nee. Mechanische energie kan afnemen terwijl totale energie wordt
        omgezet in andere vormen.

  \item \textbf{``Als de kinetische energie constant is, werken er geen
        krachten.''}\\
        Niet noodzakelijk. Bij een cirkelbeweging kan een resulterende kracht
        loodrecht op de snelheid staan en de richting veranderen zonder de
        snelheidsgrootte te veranderen.
\end{itemize}


\begin{oefening}[Misconcepties ontleden]{\oefU\ESS}

Kies drie uitspraken uit de lijst hierboven.

Voor elke uitspraak:

\begin{question}
geef een concreet tegenvoorbeeld;
\end{question}

\begin{question}
noem de relevante formule of wet;
\end{question}

\begin{question}
formuleer een fysisch correcte vervangende uitspraak.
\end{question}

\end{oefening}


% ============================================================
\FVDsection{Samenvatting}
% ============================================================

Voor een constante kracht:

\[
\boxed{
W=\vec F\cdot\Delta\vec r
=F\Delta r\cos\theta.
}
\]

Voor een veranderlijke kracht in één dimensie:

\[
\boxed{
W=\int_{x_1}^{x_2}F_x(x)\,dx.
}
\]

Grafisch is arbeid de georiënteerde oppervlakte onder de
kracht-plaatsgrafiek.

Kinetische energie:

\[
\boxed{
E_k=\frac12mv^2.
}
\]

Arbeid-energietheorema:

\[
\boxed{
W_{\mathrm{res}}=\Delta E_k.
}
\]

Voor een conservatieve kracht:

\[
\boxed{
W_{\mathrm{cons}}=-\Delta E_p.
}
\]

Nabij het aardoppervlak:

\[
\boxed{
E_{p,g}=mgh.
}
\]

Voor een ideale veer:

\[
\boxed{
E_{p,\mathrm{el}}=\frac12kx^2.
}
\]

Mechanische energie:

\[
\boxed{
E_m=E_k+E_p.
}
\]

Wanneer alleen conservatieve krachten arbeid verrichten:

\[
\boxed{
E_m=\text{constant}.
}
\]

Algemene gravitatiepotentiële energie:

\[
\boxed{
E_{p,g}(r)=-\frac{GMm}{r}.
}
\]

Voor een ideale satelliet in een cirkelbaan:

\[
\boxed{
E_m=-\frac{GMm}{2r}.
}
\]




% ============================================================
\FVDsection{Vooruitblik: van energie naar trillingen}
% ============================================================

Een massa aan een veer bezit tijdens haar beweging afwisselend kinetische en
elastische potentiële energie.

In de uiterste stand is de snelheid nul en is de elastische energie maximaal.
In de evenwichtsstand is de snelheid maximaal.

Dezelfde veer waarvoor we in dit hoofdstuk vonden:

\[
F=-kx
\]

zal in het volgende thema leiden tot:

\[
a=-\frac{k}{m}x.
\]

Daarmee verschijnt een nieuw bewegingspatroon: de harmonische trilling.

\begin{center}
\[
\boxed{
\text{arbeid en energie}
\longrightarrow
\text{trillingen}
\longrightarrow
\text{golven}.
}
\]
\end{center}

\end{document}
